\documentclass[10pt,a4paper]{amsart}
\usepackage[margin=19mm]{geometry}
\usepackage{amsmath,amssymb,mathtools}
\usepackage[scr=rsfs,cal=euler]{mathalfa}
\usepackage{xfrac}
\usepackage[unicode,hidelinks,bookmarks=false]{hyperref}
\newcommand{\ie}{i.e.\ }
\newcommand{\cf}{cf.\ }
\newcommand{\resp}{respectively\ }
\newcommand{\Subsection}{\subsection}
\providecommand{\nobreakdash}{\nobreak\hbox{-}}
\setlength{\emergencystretch}{2em}
\multlinegap0pt
\usepackage{bm}
\usepackage{bbm}
\usepackage{dsfont}
\usepackage{xspace}
\usepackage{cancel}
\usepackage{mathtools}
\usepackage{amsthm}
\makeatletter
\@ifundefined{algorithm}{\newtheorem{algorithm}{Algorithm}}{}
\@ifundefined{algorithmic}{\newtheorem{algorithmic}{Algorithmic}}{}
\@ifundefined{theorem}{\newtheorem{theorem}{Theorem}}{}
\@ifundefined{lemma}{\newtheorem{lemma}[theorem]{Lemma}}{}
\makeatother
\newcommand{\AC}{\ensuremath\mathsf{AC}}
\newcommand{\LOGSPACE}{\ensuremath\mathsf{L}}
\title[On the Complexity of Properties of Partial Bijection Semigro]{On the Complexity of Properties of Partial Bijection Semigroups}
\author{Trevor Jack}
\date{Source: arXiv:2101.00324v2 (CC BY 4.0)}
\begin{document}
\maketitle

\noindent\emph{Source and licence.} On the Complexity of Properties of Partial Bijection Semigroups. Version arXiv:2101.00324v2 (CC BY 4.0), distributed under the Creative Commons Attribution 4.0 International licence, as recorded in the arXiv licence metadata for this exact version and shown on the abstract page https://arxiv.org/abs/2101.00324v2. Full rights notice: licenses/arxiv-2101.00324-CC-BY-4.0.txt.\par\smallskip

\noindent\textit{Editorial scope.} Counted unit: the Theorem whose source label is \texttt{None} in the article (source0.tex lines 186--188, source label \texttt{None}), delivered together with the article's own complete original proof of it (source0.tex lines 189--229, one \texttt{proof} environment of 41 lines). No mathematics is written, repaired or re-derived: statement and proof are the article's own text line for line. Required context retained uncounted: lem:leftid (lines 166--169); lem:rightid (lines 176--179). Every definition, display and label of those lines is retained unchanged. Omitted from this package and not redistributed: 1--165, 170--175, 180--185, 230--470. Nothing inside a retained excerpt is omitted or altered: every retained body is delivered exactly as the article's lines are written, so no omission receipt of any kind accompanies this package. Citations: the retained excerpts carry no citation command at all, so no bibliography entry of the article is transcribed and no reference is redistributed. Reference adaptation: none was needed and none was made; every reference inside the delivered unit resolves inside it, so no unresolved reference or citation is produced. Printed slips kept literal: no printed word, symbol, hypothesis or number of the counted statement or of its proof was altered. Layout and class: the article's own class is \texttt{amsart} with packages including graphicx, algorithm, algpseudocode, xcolor, stmaryrd, amsaddr, tikz, hyperref; the wrapper is the lane's pinned \texttt{amsart} editorial wrapper at 10pt on A4, which restates the theorem-like environments the retained text needs and adds the article's own macro definitions that the retained text needs (\texttt{\textbackslash AC}, \texttt{\textbackslash LOGSPACE}), with extra wrapper packages (bm, bbm, dsfont, xspace, cancel, mathtools). The delivered numbering is the unit's own. Assets: no retained span contains a figure environment, a graphics-inclusion command, a PDF-page inclusion, a table or a TikZ picture; no image file is copied into this package, the package declares no asset dependency and no image file is redistributed with it. Licence: version arXiv:2101.00324v2 of this article is distributed under the Creative Commons Attribution 4.0 International licence, as recorded in the arXiv licence metadata for this exact version and shown on the abstract page https://arxiv.org/abs/2101.00324v2. A full rights notice is delivered at licenses/arxiv-2101.00324-CC-BY-4.0.txt. Characters: every retained excerpt is pure ASCII, so the delivered engine, XeLaTeX with its default Latin Modern text font, reports no missing character.
\begin{lemma} \label{lem:leftid}
Let $S := \langle a_1, \dots, a_k \rangle \leq T_n$. Then the idempotent power of $a_i$ is a left identity of $S$ iff
\[ \forall j \in [k] \, \forall x,y \in [n]: (x a_i = y a_i \implies x a_j = y a_j) \wedge (x a_i^2 = y a_i^2 \implies x a_i = y a_i) \]
\end{lemma}

\begin{lemma} \label{lem:rightid}
Let $S := \langle a_1, \dots, a_k \rangle \leq T_n$. Then the idempotent power of $a_i$ is a right identity of $S$ iff
\[ \forall j,\ell \in [k] \, \forall x,y \in [n]: x a_j a_i = y a_\ell a_i \implies x a_j = y a_\ell \]
\end{lemma}

\begin{theorem}
Enumerate Identities is in $\AC^0$.
\end{theorem}

\begin{proof}
For each generator $a_i$, we can define an $\AC^0$ circuit to check the formulas in Lemma~\ref{lem:leftid} and Lemma~\ref{lem:rightid}. If this circuit confirms that $a_i^\omega$ is a left or right identity, we would then like to identify the images $xa_i^\omega$ for each $x \in [n]$. Fortunately, the conditions in each of the lemmas offer us a way around direct computation.

Consider any $x,y \in [n]$. Since $xa_i = xa_i^{\omega+1}$, then $xa_i = ya_i^2$ implies $xa_i^{\omega+1} = ya_i^2$. By the conditions in each Lemma, this implies that $xa_i^\omega = ya_i$. So, the value of $xa_i^\omega$ equals $ya_i$ where $y$ satisfies $xa_i = ya_i^2$. Then, for each $x \in [n]$, each generator $a_i$ that satisfies the formulas in either Lemm~\ref{lem:leftid} or Lemma~\ref{lem:rightid}, and each $y \in [n]$, we define an $\AC^0$ to check if $xa_i = ya_i^2$. If the circuit accepts, then $xa_i^\omega = ya_i$

\iffalse
We claim that Algorithm~\ref{alg:identities} solves Enumerate Identities deterministically using $O(\log(k)+\log(n))$ space.

\emph{Correctness:} Lines 4-9 check which generators $a_1,\dots,a_k$ satisfy the conditions of Lemma~\ref{lem:leftid} or Lemma~\ref{lem:rightid}. If $a_i$ satisfies either of these conditions, we then need to compute the images of the idempotent power $a_i^\omega$. Doing so directly could require more than logarithmic space, but fortunately the conditions in each of the lemmas offer us a way around direct computation.

For any $x,y \in [n]$, each of the conditions guarantees that $xa_i^{\omega+1} = y a_i^2$ implies $xa_i^\omega = ya_i$. Since $xa_i = xa_i^{\omega+1}$, then $xa_i = ya_i^2$ implies $xa_i^\omega = ya_i$. Thus, the value of $xa_i^\omega$ equals $ya_i$ where $y$ satisfies $xa_i = ya_i^2$. Lines 10-17 iterate through each $x \in [n]$ to find a point $y \in [n]$ satisfying $xa_i = ya_i^2$ and then prints $ya_i$ as the value for $xa_i^\omega$.

\emph{Runspace:} Lines 1-3 use $\log(n)+2$ space. Line 4 uses $2(\log(k)+\log(n))$ space. Lines 5-8 use $\log(n)$ space for each of the 12 points in the three clauses. Each clause can write over the space used by the previous clause, so Lines 4-9 use $2\log(k)+6\log(n)$ space. Lines 10-16 only ever need $4\log(n)$ space and can write over the space used by Lines 5-8. Thus, the whole algorithm uses $O(\log(k)+\log(n))$ space.
  \begin{algorithm}
  \caption{$\LOGSPACE$ algorithm for Enumerate Identities}
    \label{alg:identities}
    \begin{algorithmic}[1]
      \Input{$a_1,\dotsc,a_k \in T_n$}
      \Output{Left and right identities of $\langle a_1,\dots,a_k\rangle$}
      \ForAll{$i \in [k]$}
          \State isLeftId := true
          \State isRightId := true
          \ForAll{$j,\ell \in [k]$ and $x,y \in [n]$}
             \If{$(xa_i = ya_i \land xa_j \neq ya_j) \lor (xa_i^2 = ya_i^2 \land xa_i \neq ya_i)$}
                \State isLeftId := false
             \EndIf
             \LineIf{$(xa_ja_i = ya_\ell a_i \land xa_j \neq ya_\ell)$}{isRightId := false}
          \EndFor
          \If{isLeftId or isRightId}
            \ForAll{$x,y \in [n]$}
                \If{$xa_i = ya_i^2$}
                    \State Print $ya_i$
                    \State Skip to next value for $x$
                \EndIf
            \EndFor
          \EndIf
      \EndFor
    \end{algorithmic}
  \end{algorithm}
  \fi
\end{proof}

\end{document}
